\scanpage{10}
on the other hand that $P$ is locally of finite presentation over $k$, it commutes (as afunctor) to filtering direct limits of rings, and therefore it is known when we know it's restriction to arguments which are of finite type over $k$, and a fortiori \deleted{\unsure{are}} noetherian. For such \inserted{\handwritten{(i.e. noetherian)}} arguments, using the existence theorem for coherent sheaves, we can give a more geometric expression for the right hand side of (18), \deleted{namely} Assume for simplicity $T=\operatorname{Spec}(B)$ affine, and let
\[
C=A\handwritten{\widehat{\otimes}}_k B=\lim A_n\handwritten{\otimes}_k B\ ,
\]
which is an adic noetherian ring augmented to $B$ (the augmentation ideal being an ideal of definition). Consider $\deleted{\unsure{X_T|Y}}\inserted{\handwritten{X\widehat{\otimes}_k B}}=X\handwritten{\otimes}_A C$\inserted{\deleted{\handwritten{$\unsure{=X\widehat{\otimes}_k B}$}}}, which is a proper scheme over $C$. We then have
\[
\operatorname{Pic}(X\handwritten{\widehat{\otimes}}_k B)\handwritten{\xrightarrow{\sim}}\underset{\handwritten{n}}{\underset{\handwritten{\longleftarrow}}{\lim}}\operatorname{Pic}(X_{nB}),\deleted{\unsure{\longrightarrow P(B)}}\handwritten{\operatorname{Pic}(X\widehat{\otimes}_{\unsure{n|k}}B)/\operatorname{Pic}(B)\hookrightarrow P(B)}
\tag*{(19)}
\]
and thus on \uline{noetherian} arguments $B$, $P$ appears as the étale sheaf associated to the functor $B\handwritten{\mapsto}\operatorname{Pic}(X\handwritten{\widehat{\otimes}}_k B)\handwritten{/\operatorname{Pic}(B)}$ \inserted{\handwritten{which is a \uline{separated} presheaf,}} this statement being complemented by the description given above of the group for obstructions for an element of $P(B)$ to \unsure{beleng|belong} to $\operatorname{Pic}(X\handwritten{\widehat{\otimes}}_k B)$.

As for $\deleted{\unsure{Q(B)}}B\handwritten{\mapsto}Q(B)$, this appears as the \deleted{\unsure{sheaf}} étale sheaf associated with the functor
\[
B\handwritten{\longmapsto}\operatorname{Pic}(X\handwritten{\widehat{\otimes}}_k B)/\inserted{\handwritten{(\operatorname{Pic}(B)+}}\underline{D}(B))\ (\handwritten{\unsure{\hookrightarrow|\longleftrightarrow}}Q(B))
\tag*{(20)}
\]
where $\underline{D}$ is the twisted constant group over $k$ defined above. For an element of $Q(B)$ to belong to the first hand side there are two successive obstructions, the first is in $H^1(\operatorname{Spec}(B),\underline{D})$, which vanishes \inserted{\handwritten{for instance}} if $\underline{D}$ is constant (for instance $k$ separably closed) and $B$ normal, the second is in the Brauer group of $B$, and vanishes \inserted{\handwritten{for instance}} if $X_o$ has a zero-cycle of degree $1$.

Assume for simplicity that the irreducible components of $X_o$ are
% archive sheet number: 9
